\section{The radial envelope and the norm certificate}
\label{bndnorm:section}\label{energy:section}
The same scalar comparison gives the norm, gap and joint certificates.
Put $\mathcal A(v)=2H_0v+H_1v^2$. For $P>q\ge0$, define
\begin{equation}\label{boundarygap:time}
 \mathcal T(U;P,q)=\int_0^U
 \frac{dv}{\sqrt{\min\{q^2+\mathcal A(v),
                  P^2+\mathcal A(U)-\mathcal A(v)\}}}.
\end{equation}
For $P>0$ and $0<r\le P/m$, let $A(P,r)$ solve
\begin{equation}\label{compact:radial-inverse}
 \mathcal T(A(P,r);P,P-mr)=r,
 \qquad A(P,0)=0.
\end{equation}

\begin{lemma}[Radial comparison and sharpness]\label{compact:radial}
The inverse in \eqref{compact:radial-inverse} exists and is unique.
Let $u\in C^1([0,r])$ have an absolutely continuous derivative and satisfy
\[
 u(0)=0,\quad u'(0)\ge0,\quad
 u'(s)\ge u'(0)+ms,\quad u'(r)\le P,
 \quad |u''(s)|\le H_0+H_1u(s)\quad\text{a.e.}
\]
Then $u(r)\le A(P,r)$. Equality holds for a trajectory whose curvature
is first $H_0+H_1u$ and then $-(H_0+H_1u)$, with initial slope
$P-mr$ and terminal slope $P$. Smooth trajectories with these exact
slopes and the same constants approach its value.
For fixed $P$, $A(P,r)$ is continuous and strictly increasing in $r$.

For $H_1=0$,
\begin{equation}\label{compact:radial-constant}
 A(P,r)=Pr+\beta r^2,\qquad
 \beta=\frac{H_0^2-2mH_0-m^2}{4H_0}.
\end{equation}
For $H_1>0$, put $k=\sqrt{H_1}$, $a=H_0/H_1$, $T=kr$, $q=P-mr$,
and, for $0\le b\le\min\{T,\pi/4\}$, set
\[
 Y_b=a\cosh(T-b)+(q/k)\sinh(T-b),\qquad
 Q_b=ak\sinh(T-b)+q\cosh(T-b).
\]
The unique root of
\begin{equation}\label{compact:radial-switch}
 Q_b\cos b-kY_b\sin b=P
\end{equation}
is interior, and
\begin{equation}\label{compact:radial-value}
 A(P,r)=Y_b\cos b+(Q_b/k)\sin b-a.
\end{equation}
\end{lemma}

\begin{proof}
Write $p=u'$, $U=u(r)$ and $q=P-mr$. The assumptions imply
$0\le p(0)\le q$ and $p(s)>0$ for $s>0$. The absolutely continuous
chain rule applied to $p^2-\mathcal A(u)$ and $p^2+\mathcal A(u)$ gives
\begin{equation}\label{energy:two-sided}
 p(s)^2\le
 \min\{q^2+\mathcal A(u(s)),P^2+\mathcal A(U)-\mathcal A(u(s))\}.
\end{equation}
Changing variables from $s$ to $u(s)$ yields
$\mathcal T(U;P,q)\le r$. When $q=0$, the only endpoint singularity
is of order $v^{-1/2}$ and is integrable. This argument also covers
joins at which the second derivative jumps.

We verify strict inversion. If $\mathcal A(U)\le P^2-q^2$, the first
expression in the minimum is active throughout, so $\mathcal T$ is
strictly increasing. Otherwise the crossing $v_c\in(0,U)$ satisfies
$2\mathcal A(v_c)=P^2+\mathcal A(U)-q^2$.
For $H_1>0$, set
\[
 z=U+a,\qquad R^2=z^2+P^2/k^2,\qquad
 w^2=\tfrac12(R^2+a^2-q^2/k^2),\qquad w=v_c+a.
\]
Then
\[
 \mathcal T=
 \int_a^w\frac{dy}{\sqrt{k^2y^2+q^2-k^2a^2}}
 +\frac1k\{\arcsin(z/R)-\arcsin(w/R)\}.
\]
Differentiation cancels the crossing terms and gives
\begin{equation}\label{boundarygap:time-derivative}
 \mathcal T_U=\frac{P}{k^2R^2}
       +\frac{zw}{kR^2\sqrt{R^2-w^2}}>0.
\end{equation}
The formulas agree at the branch boundary. The first integral diverges
as $U\to\infty$, while the second is nonnegative. For $H_1=0$,
\[
 \mathcal T(U;P,q)=
 \begin{cases}
 (\sqrt{q^2+2H_0U}-q)/H_0,&2H_0U\le P^2-q^2,\\[1mm]
 (\sqrt{2q^2+2P^2+4H_0U}-q-P)/H_0,&2H_0U>P^2-q^2.
 \end{cases}
\]
Thus in both cases $\mathcal T$ is continuous, strictly increasing
from zero to infinity. This proves the inverse and the upper bound.

For $H_1>0$, solve $u''=H_0+H_1u$ from $u(0)=0$, $u'(0)=q$ for
time $(T-b)/k$, then solve $u''=-(H_0+H_1u)$ for time $b/k$.
Its terminal slope is $F(b)=Q_b\cos b-kY_b\sin b$, with
$F'(b)=-2kY_b\cos b<0$. At $b=0$,
$F(0)=ak\sinh T+q\cosh T>q+H_0r\ge P$.
At the other endpoint it is either $q\cos T-ka\sin T<P$, or
$(q-ka)e^{-(T-\pi/4)}/\sqrt2<P$.
Hence \eqref{compact:radial-switch} has exactly one interior root.
On its first arc the derivative is at least $q+ms$; on its second
arc it decreases to $P=q+mr$, so the same inequality holds.
The loss stays nonnegative, both energy bounds are attained on their
respective arcs, and \eqref{compact:radial-value} follows.
For $H_1=0$, switching from $H_0$ to $-H_0$ at
$r(H_0+m)/(2H_0)$ gives \eqref{compact:radial-constant}.
If also $H_0=m$, this trajectory is the smooth function $qs+ms^2/2$.

In every other case the switch is interior and its slope exceeds $P$.
Replace its sign by a smooth nonincreasing transition in $[-1,1]$
over an interval of width $2\epsilon$, and solve
$u''=\chi_{\epsilon,c}(H_0+H_1u)$ with the same initial data.
Continuous dependence of this linear ODE gives uniform convergence
of $u,u'$ as $\epsilon\downarrow0$. The limiting terminal slope has
strictly positive derivative with respect to the physical switch $c$:
$2(H_0+H_1u(c))\cos(k(r-c))$, or $2H_0$ if $H_1=0$.
The intermediate value theorem therefore permits a nearby transition
center for which $u'(r)=P$ exactly. Before the transition
$u'(s)\ge q+ms$; throughout the transition it stays above $P$ by a
fixed positive margin; afterwards it decreases to $P$. The same RSI
and curvature bounds hold, with unchanged constants and endpoint slopes.
The solutions extend smoothly to a neighborhood of the interval and
their losses converge to $A(P,r)$.

Finally let $0<r_1<r_2\le P/m$, $d=r_2-r_1$, and $q_i=P-mr_i$.
Prefix a near-extremizer at $r_1$ by $q_2t+mt^2/2$ on $[0,d]$.
Its final slope is $q_1$ and its value is
$a_0=q_2d+md^2/2>0$. Continuing with $a_0+u(t-d)$ preserves the
derivative lower bound and can only enlarge the permitted curvature.
The joined derivative is absolutely continuous, so the energy comparison
already proved gives
\begin{equation}\label{compact:radial-growth}
 A(P,r_2)\ge A(P,r_1)+q_2d+md^2/2>A(P,r_1).
\end{equation}
Continuity for $r>0$ follows from the strict switch, and at zero from
the all-positive-curvature bound, which tends to zero. These comparison
trajectories do not enlarge the feasible set of the original problem.
\end{proof}

Define $\mathcal E(P)=A(P,P/m)$ for $P>0$ and $\mathcal E(0)=0$.
The preceding formulas evaluate it with one scalar root when $H_1>0$;
for $H_1=0$ they give
\begin{equation}\label{energy:constant}
 \mathcal E(P)=
 \frac{P^2(H_0^2+2mH_0-m^2)}{4m^2H_0}.
\end{equation}

\begin{theorem}[Norm certificate under shifted RSI]
\label{bndnorm:main}\label{energy:main}
Under \eqref{boundarygap:assumptions}, every source satisfies
\begin{equation}\label{bndnorm:observable}
 \delta(x)\le\mathcal E(P(x)).
\end{equation}
For $x\ne x^*$, with radial slopes $q=u'(0)$ and $p=u'(r)$,
\begin{equation}\label{bndnorm:radial-deficit}
 \delta(x)+\mathcal E(q)\le\mathcal E(p)\le\mathcal E(P(x)).
\end{equation}
The function $\mathcal E$ is continuous and, for $0\le P_1<P_2$,
\begin{equation}\label{compact:norm-growth}
 \mathcal E(P_2)\ge\mathcal E(P_1)+(P_2^2-P_1^2)/(2m).
\end{equation}
For each $P>0$, the bound is sharp with $P(x)=P$ exactly, including
among smooth interval examples with $f'(x^*)>0$.
\end{theorem}
\begin{proof}
The radial reduction and Lemma~\ref{compact:radial} give
$\delta(x)\le A(P(x),r)\le A(P(x),P(x)/m)$.
Extend a stationary extremal trajectory at $P_1$ from $P_1/m$ to
$P_2/m$ with derivative $ms$. The join has an absolutely continuous
derivative and curvature $m\le H_0+H_1u$ on the extension.
The same comparison proves \eqref{compact:norm-growth}, including
$P_1=0$ by the quadratic trajectory. Continuity follows from the
switch formula; at zero use the all-positive-curvature upper bound.

For \eqref{bndnorm:radial-deficit}, prefix the original ray by smooth
stationary trajectories on $[0,q/m]$ with terminal slope $q$ and losses
$a_j\to\mathcal E(q)$. Continue with $a_j+u(t-q/m)$.
The first derivatives match, shifted RSI becomes $v'(t)\ge mt$,
and the nonnegative loss offset preserves the curvature bound.
The joined ray has terminal slope $p$ and length at most $p/m$.
The comparison gives $a_j+\delta(x)\le\mathcal E(p)$; pass to the
limit. The case $q=0$ uses an empty prefix.

Sharpness with a stationary minimum is Lemma~\ref{compact:radial}
at $r=P/m$. To make the optimal derivative positive, take
$q_j\downarrow0$, use stationary near-extremizers $v_j$ at endpoint
slope $P-q_j$ on $[0,(P-q_j)/m]$, and set $f_j(s)=v_j(s)+q_js$.
The constants are unchanged, $f_j'(0)=q_j>0$, $f_j'(r_j)=P$, and
$f_j(r_j)\to\mathcal E(P)$ by continuity.
\end{proof}

\begin{proposition}[A fixed initial slope]\label{joint:fixed-slope}
For $P>q\ge0$, the sharp radial loss with initial slope $q$ and terminal
slope $P$ is
\begin{equation}\label{joint:fixed-slope-time}
 \mathcal T(\mathcal E_q(P);P,q)=(P-q)/m,
 \qquad \mathcal E_q(P)\le\mathcal E(P)-\mathcal E(q).
\end{equation}
Smooth interval functions with both slopes and all constants fixed
approach this value. If $H_1=0$, then
\begin{equation}\label{joint:fixed-slope-constant}
 \mathcal E_q(P)=\frac{P^2-q^2}{2m}
       +\frac{H_0^2-m^2}{4H_0m^2}(P-q)^2,
\end{equation}
and the deficit bound has excess
\[
 \mathcal E(P)-\mathcal E(q)-\mathcal E_q(P)
 =\frac{q(P-q)(H_0^2-m^2)}{2H_0m^2}.
\]
\end{proposition}
\begin{proof}
RSI gives $r\le(P-q)/m$, and the energy comparison with these fixed
slopes gives $\mathcal T(u(r);P,q)\le r$. Strict inversion proves the
bound. Lemma~\ref{compact:radial} at $r=(P-q)/m$ proves smooth sharpness;
\eqref{bndnorm:radial-deficit} proves the deficit. Substituting this
radius in \eqref{compact:radial-constant} gives the two algebraic formulas.
\end{proof}

\begin{corollary}[The stationary full-gap certificate]\label{energy:gap}
If $\nabla f(x^*)=0$, set $W(g)=\mathcal E(\sqrt{mg})$. Then
\begin{equation}\label{energy:gap-certificate}
 \delta(x)\le W\!\left(\min\{G(x),P(x)^2/m\}\right).
\end{equation}
The function $W$ is continuous, strictly increasing, unbounded, and
$W(g)\ge g/2$. The bound is sharp when the gap is fixed exactly.
\end{corollary}
\begin{proof}
For fixed $g>0$, define $U(r)$ by
$\mathcal T(U(r);g/r,0)=r$. Write $P=g/r$ and let $v_c$ be the
crossing of the two energy bounds. In the mixed branch,
\[
 \mathcal T+P\mathcal T_P
 =\int_0^{v_c}\frac{dv}{\sqrt{\mathcal A(v)}}
 +\int_{v_c}^{U}
   \frac{\mathcal A(U)-\mathcal A(v)}
        {[P^2+\mathcal A(U)-\mathcal A(v)]^{3/2}}\,dv>0.
\]
In the forward-only branch $\mathcal T_P=0$. Thus
$U'(r)=(\mathcal T+P\mathcal T_P)/(r\mathcal T_U)>0$.
Since $r\le\sqrt{g/m}$, the stationary radial comparison gives
$\delta(x)\le U(\sqrt{g/m})=W(g)$. Combine this with the norm
bound. The properties of $W$ follow from \eqref{compact:norm-growth}.
The stationary interval examples have $G=P^2/m$ exactly and prove
sharpness. This argument uses stationarity; the shifted gap envelope
in the next section is different.
\end{proof}
